\section*{ОБОЗНАЧЕНИЯ И СОКРАЩЕНИЯ}

ВКР -- выпускная квалификационная работа.

ОПОП ВО -- основная профессиональная образовательная программа высшего образования.

ПО -- программное обеспечение.

ЮЗГУ -- Юго-Западный государственный университет.

AABB (Axis-Aligned Bounding Box) -- ограничивающий параллелепипед, выровненный по осям координат.

API (Application Programming Interface) -- программный интерфейс приложения.

EBO (Element Buffer Object) -- объект буфера индексов OpenGL.

ECS (Entity-Component-System) -- архитектурный подход к построению программных систем, основанный на сущностях, компонентах и системах обработки.

GLFW -- библиотека, используемая для создания окон, управления контекстом OpenGL и обработки пользовательского ввода.

GLSL (OpenGL Shading Language) -- язык программирования шейдеров OpenGL.

GPU (Graphics Processing Unit) -- графический процессор.

JAR (Java Archive) -- архивный формат, используемый для распространения Java-классов и связанных с ними ресурсов.

JDK (Java Development Kit) -- комплект средств разработки для платформы Java.

JVM (Java Virtual Machine) -- виртуальная машина Java.

LWJGL (Lightweight Java Game Library) -- библиотека, предоставляющая Java-приложениям доступ к низкоуровневым графическим и системным API.

OpenGL (Open Graphics Library) -- кроссплатформенный графический API для вывода двухмерной и трехмерной графики.

SAT (Separating Axis Theorem) -- теорема разделяющей оси, используемая для проверки пересечения выпуклых тел.

SDK (Software Development Kit) -- комплект средств разработки.

STB -- набор вспомогательных библиотек, применяемых в том числе для загрузки изображений и других мультимедийных ресурсов.

UML (Unified Modeling Language) -- язык графического описания для моделирования программных систем.

VAO (Vertex Array Object) -- объект массива вершин OpenGL.

VBO (Vertex Buffer Object) -- объект буфера вершин OpenGL.
